\subsection{Нецентральные удары}
%%%%%%%%%%%%%%%%%%%%%%%%%%%%%%%%%%%%%%%%%%%%%%%%%%%%
%%%%%%%%%%%%%%%%%%%%%%%%%%%%%%%%%%%%%%%%%%%%%%%%%%%% 
\z[{\begin{wrapfigure}[5]{r}{0.2\textwidth}%
		\vspace{-15pt}%
		\begin{flushright}%
 			\includegraphics[width=0.2\textwidth]{7.9.1}%
		\end{flushright}%
	\end{wrapfigure}%
	}]%
	{Частица массы $m_1$ налетает на шар массы $m_2$. Направление её движения составляет угол $\alpha$ с нормалью к поверхности шара. Под каким углом к этой нормали отскочит от шара частица, если шар сначала покоился, а удар упругий?}
%
{$\tg \beta = \tg \alpha \cdot \frac{m_1 + m_2}{m_2 - m_1}$}
%%%%%%%%%%%%%%%%%%%%%%%%%%%%%%%%%%%%%%%%%%%%%%%%%%%%
\z{Тяжёлая частица массы $m_1$ сталкивается с покоящейся лёгкой частицей массы $m_2$. На какой наибольший угол может отклониться тяжёлая частица в результате упругого удара?}
%
{$\sin \alpha = \frac{m_2}{m_1}$}
%%%%%%%%%%%%%%%%%%%%%%%%%%%%%%%%%%%%%%%%%%%%%%%%%%%%	
\z[{\begin{wrapfigure}[5]{r}{0.3\textwidth}%
		\vspace{-15pt}%
		\begin{flushright}%
 			\includegraphics[width=0.3\textwidth]{7.9.2}%
		\end{flushright}%
	\end{wrapfigure}%
	}]%
	{При упругом столкновении налетающей частицы с покоящейся первая полетела под углом $\alpha$ к направлению первоначального движения, а вторая -- под углом $\beta$. Найдите отношение масс этих частиц.}
%
{$\frac{m_1}{m_2} = \frac{\sin^2(\alpha + \beta) - \sin^2 \beta}{\sin^2 \alpha}$}

%%%%%%%%%%%%%%%%%%%%%%%%%%%%%%%%%%%%%%%%%%%%%%%%%%%% 
\z{Частица массы $m_1$ налетела со скоростью $v$ на неподвижную частицу массы $m_2$, которая после упругого удара полетела под углом $\alpha$ к первоначальному направлению движения налетающей частицы. Определите скорость частицы массы $m_2$ после удара.}
%
{$v_2 = \frac{2v_{10} m_1}{m_2 + m_1} \cos \alpha$}
%%%%%%%%%%%%%%%%%%%%%%%%%%%%%%%%%%%%%%%%%%%%%%%%%%%%	
\z[{\begin{wrapfigure}[5]{r}{0.3\textwidth}%
		\vspace{-10pt}%
		\begin{flushright}%
 			\includegraphics[width=0.3\textwidth]{7.9.3}%
		\end{flushright}%
	\end{wrapfigure}%
	}]%
	{На массивный гладкий цилиндр радиусом $R$, движущийся поступательно со скоростью $u$, налетает маленький шарик, движущийся навстречу цилиндру перпендикулярно его оси со скоростью $v$ (см. рисунок). Расстояние между линией, вдоль которой движется шарик, и плоскостью, в которой движется ось цилиндра, равно $L$ ($L < R$). Найдите величину скорости шарика $v_1$ после абсолютно упругого удара о цилиндр. Сила тяжести отсутствует.}
%
{$v_1 = \sqrt{v^2 + 4u(u + v)\left(1 - \frac{L^2}{R^2}\right)}$}


